% !TeX root = ../tesis.tex
%%%%%%%%%%%%%%%%%%%%%%%%%%%%%%%%%%%%%%%%%%%%%%%%%%%%%%%%%%%%%%%%%%%
% SECCIÓN 5.2-3 — RESULTADOS VFT: FASE 3
% Análisis de red: tiempo de viaje, centralidad y vulnerabilidad
%%%%%%%%%%%%%%%%%%%%%%%%%%%%%%%%%%%%%%%%%%%%%%%%%%%%%%%%%%%%%%%%%%%
\section{\bajoRevision{Fase 3: Análisis de Red y Vulnerabilidad}}
\label{sec:resultados-fase3}

\notaRevision{Los tres indicadores de esta fase —Tiempo de Viaje Promedio,
Centralidad de Intermediación y Vulnerabilidad Estructural— requieren que las
Fases~1 y~2 estén completas. A la fecha de cierre del documento, los algoritmos
de caminos mínimos no han sido ejecutados sobre la versión final ponderada del
grafo. Las fórmulas, interpretaciones esperadas y estructura de resultados
anticipados se presentan como fundamento para la etapa de cálculo.}

La Fase~3 representa la capa de análisis de mayor complejidad computacional del
Modelo VFT. A diferencia de los indicadores de las fases anteriores, que operan
sobre la geometría estática de la red o sobre los pesos individuales de aristas y
nodos, los tres indicadores de esta fase requieren la ejecución de \textbf{algoritmos
de caminos mínimos} sobre la totalidad del grafo ponderado, lo que implica tiempos
de cómputo del orden de horas para un grafo de 11{,}115 nodos y 45{,}679 aristas.

La secuencia de dependencias es estricta: el \textbf{Tiempo de Viaje Promedio}
($T$) exige el grafo con todos sus pesos dinámicos de Fase~2; la
\textbf{Centralidad de Intermediación} ($B$) requiere los caminos mínimos
calculados en $T$; y el análisis de \textbf{Vulnerabilidad Estructural}
($\Delta E$) necesita $B$ para identificar los nodos críticos sobre los cuales
simular fallas.

% ----------------------------------------------------------------
\subsection{\bajoRevision{Tiempo de Viaje Promedio ($T$)}}
\label{subsec:tiempo-viaje-promedio}
% ----------------------------------------------------------------

\notaRevision{Pendiente de cálculo. Requiere Fases 1 y 2 completas más
ejecución del algoritmo \dijkstra{} sobre el grafo final ponderado.}

El \textbf{Tiempo de Viaje Promedio} ($T$) mide la eficiencia global de la red
como sistema: cuánto tarda en promedio un pasajero en desplazarse entre cualquier
par de estaciones alcanzables. Es la métrica de referencia que sintetiza, en un
solo valor, el efecto combinado de la geometría de la red, los tiempos de traslado
con fricción y los costos de transbordo.

\subsubsection{Fórmula}
\label{subsubsec:formula-t}

\begin{equation}
    T = \frac{1}{N(N-1)} \sum_{\substack{i,j \in \mathcal{V} \\ i \neq j}} t(i,j)
    \label{eq:tiempo-viaje-promedio}
\end{equation}

\begin{itemize}
    \item $\mathcal{V}$: conjunto de los $N = 11{,}115$ nodos topológicos del
    grafo.
    \item $t(i,j)$: tiempo de viaje mínimo del nodo $i$ al nodo $j$ en minutos,
    calculado mediante el algoritmo \dijkstra{} sobre el grafo ponderado con los
    atributos de Fase~2.
    \item El denominador $N(N-1)$ normaliza sobre todos los pares ordenados de
    nodos alcanzables. Los pares sin ruta (componentes desconectadas) se excluyen
    del promedio.
\end{itemize}

\subsubsection{Resultado Esperado e Interpretación}
\label{subsubsec:t-resultado-esperado}

El indicador $T$ caracteriza la red en su conjunto como sistema físico: un valor
bajo indica una red eficiente donde cualquier par de estaciones está bien
conectado en el plano temporal; un valor alto sugiere desconexiones estructurales
o cuellos de botella en los transbordos. La comparación de $T$ antes y después de
la remoción simulada de nodos críticos (Sección~\ref{subsec:vulnerabilidad}) es
el mecanismo central del análisis de vulnerabilidad.

Adicionalmente, el modelo contempla una variante ponderada por demanda
$T_{\text{dem}}$, calculable una vez que el \termino{Node Score} de Fase~2 esté
disponible, donde los pares O-D con mayor afluencia en ambos extremos contribuyen
más al promedio. La diferencia $T_{\text{dem}} - T$ cuantificaría el sesgo de la
red: si las rutas más demandadas son sistemáticamente más lentas que el promedio
estructural, ese diferencial tiene implicaciones directas para la política de
mejora de infraestructura.

% ----------------------------------------------------------------
\subsection{\bajoRevision{Centralidad de Intermediación ($B$)}}
\label{subsec:res-centralidad-intermediacion}
% ----------------------------------------------------------------

\notaRevision{Pendiente de cálculo. Requiere la matriz de caminos mínimos
producida por el indicador $T$ (Sección~\ref{subsec:tiempo-viaje-promedio}).}

La \textbf{Centralidad de Intermediación} ($B$) mide cuántas veces cada nodo
aparece en los caminos más cortos entre todos los pares de estaciones de la red.
Un nodo con $B$ alto es un \termino{hub} de paso: su eliminación o degradación
afecta desproporcionadamente al número de rutas disponibles para los pasajeros.
En la hipótesis central de esta investigación, los corredores anulares
—\textit{Periférico}, Circuito Interior— concentran nodos de alta intermediación
que articulan los flujos radiales de la red sin ser ellos mismos estaciones
terminales de alto volumen.

\subsubsection{Fórmula}
\label{subsubsec:formula-b}

\begin{equation}
    B(v) = \frac{1}{(N-1)(N-2)} \sum_{\substack{s,t \in \mathcal{V} \\ s \neq v \neq t}}
    \frac{\sigma_{st}(v)}{\sigma_{st}}
    \label{eq:centralidad-intermediacion}
\end{equation}

\begin{itemize}
    \item $\sigma_{st}$: número total de caminos más cortos (en tiempo) entre
    los nodos $s$ y $t$ en el grafo ponderado.
    \item $\sigma_{st}(v)$: número de esos caminos que pasan por el nodo $v$.
    \item El factor de normalización $1/[(N-1)(N-2)]$ escala el valor al rango
    $[0, 1]$ para un grafo de $N$ nodos.
\end{itemize}

El cálculo se implementa mediante la función \texttt{betweenness\_centrality()}
de \textsc{NetworkX} con el atributo \texttt{weight} como peso de las aristas,
de modo que los caminos mínimos se definen en tiempo, no en número de saltos.

\subsubsection{Resultado Esperado e Interpretación}
\label{subsubsec:b-resultado-esperado}

Se espera que la distribución de $B$ sobre los 11{,}115 nodos sea
\termino{heavy-tailed}: un pequeño número de nodos concentrará valores altos de
intermediación (nodos de paso obligado), mientras que la mayoría de estaciones
tendrá valores cercanos a cero. Este patrón es consistente con la distribución
de red libre de escala (\termino{scale-free}) detectada en la Fuerza Capilar
(Fase~1).

La hipótesis de investigación predice que al menos algunos de los nodos de
mayor $B$ se ubicarán sobre ejes circulares o de transferencia que no son
estaciones de origen-destino frecuente, sino \termino{hubs} de paso que articulan
rutas radiales. Identificar estos nodos es el prerrequisito para el análisis de
vulnerabilidad de la Sección~\ref{subsec:vulnerabilidad}.

% ----------------------------------------------------------------
\subsection{\bajoRevision{Vulnerabilidad Estructural ($\Delta E$)}}
\label{subsec:vulnerabilidad}
% ----------------------------------------------------------------

\notaRevision{Pendiente de cálculo. Requiere la Centralidad de Intermediación
$B$ (Sección~\ref{subsec:res-centralidad-intermediacion}) para seleccionar los nodos
críticos sobre los cuales ejecutar las simulaciones de falla.}

El análisis de \textbf{Vulnerabilidad Estructural} ($\Delta E$) cuantifica en
qué medida la red pierde eficiencia global cuando uno o varios nodos críticos
son removidos. El indicador mide el impacto de fallas selectivas, no aleatorias:
los nodos candidatos a la remoción son aquellos con mayor centralidad de
intermediación $B$, precisamente porque son los que mayor impacto tienen sobre
la conectividad del sistema.

\subsubsection{Fórmula}
\label{subsubsec:formula-deltae}

\begin{equation}
    \Delta E = \frac{E_0 - E_f}{E_0}
    \label{eq:vulnerabilidad}
\end{equation}

\begin{itemize}
    \item $E_0$: eficiencia global de la red original, definida como el inverso
    del tiempo de viaje promedio: $E_0 = 1/T$.
    \item $E_f$: eficiencia global del grafo residual tras la remoción del nodo
    (o conjunto de nodos) crítico.
    \item $\Delta E \in [0, 1]$: un valor cercano a 1 indica que la remoción
    del nodo colapsa la eficiencia de la red; un valor cercano a 0 indica que la
    red es robusta ante esa falla específica.
\end{itemize}

\subsubsection{Protocolo de Simulación}
\label{subsubsec:protocolo-simulacion}

El análisis se ejecutará en tres escalas de falla:

\begin{enumerate}
    \item \textbf{Falla puntual:} remoción de un único nodo (el de mayor $B$)
    y recálculo de $T_f$. Establece el impacto del nodo más crítico individual.
    \item \textbf{Falla por corredor:} remoción simultánea de todos los nodos
    de un eje identificado (e.g., todos los nodos del corredor del Periférico
    Oriente). Mide la vulnerabilidad sistémica del eje completo.
    \item \textbf{Falla en cascada:} remoción secuencial de los $k$ nodos de
    mayor $B$ residual (recalculado tras cada remoción). Simula el escenario de
    falla progresiva en eventos de alta disrupción.
\end{enumerate}

\subsubsection{Resultado Esperado e Interpretación}
\label{subsubsec:deltae-resultado-esperado}

La hipótesis central de esta investigación plantea que la red de transporte de
la \zmvm presenta \textbf{vulnerabilidad asimétrica}: nodos periféricos de alta
intermediación —ubicados en corredores anulares donde convergen rutas radiales—
producen un $\Delta E$ desproporcionadamente mayor al esperado por su posición
geográfica, en comparación con estaciones terminales de alta afluencia como
Pantitlán o Indios Verdes.

Si la simulación confirma esta hipótesis, el análisis proporciona una base
cuantitativa para priorizar inversión en infraestructura de respaldo (rutas
alternativas, capacidad de transbordo) en nodos de alta intermediación pero baja
visibilidad operativa. Si la refuta, el resultado es igualmente valioso: indicaría
que la red posee redundancia suficiente en sus ejes circulares, y que la
vulnerabilidad real se concentra en las terminales radiales de mayor afluencia.
